% !PRINTED 113
\weberpage{113}

Hiernach findet man, wenn man in (2) die Formeln § 31, (6), (11)
anwendet, für die linearen Fundamentaltransformationen von
$\eta(\omega)$
\weq{4}{\eta(\omega+1) = e^{\frac{\pi i}{12}}\eta(\omega),}
\weq{5}{\eta\left(-\frac{1}{\omega}\right) = \sqrt{-i\omega}\,\eta(\omega),}
von denen die erste auch unmittelbar aus der Reihendarstellung
(3) folgt, während die zweite nicht so leicht durch direkte
Umformung der Reihen bewiesen werden kann.

Wir gehen über zur Betrachtung der beiden Haupttransformationen
zweiter Ordnung der $\eta$-Funktion.

Aus § 32, (11) erhält man durch Differentiation und Nullsetzen
von $u$
\[ 2\vartheta_{01}(0, 2\omega)\eta(2\omega)^3 = \vartheta_{10}\eta(\omega)^3, \]
also wenn man ins Quadrat erhebt und § 32, (8) anwendet:
\[ 4\vartheta_{00}\vartheta_{10}\vartheta_{01}\eta(2\omega)^6 = \vartheta_{10}^3\eta(\omega)^6, \]
und mit Anwendung von (2) und Ausziehen der Kubikwurzel:
\weq{6}{2\eta(2\omega)^2 = \vartheta_{10}\eta(\omega).}

Ersetzt man in (6) $\omega$ durch $\dfrac{\omega}{2}$, so folgt
\[ 2\eta(\omega)^2 = \vartheta_{10}\left(0, \frac{\omega}{2}\right)\eta\left(\frac{\omega}{2}\right), \]
und wenn man quadriert und die Formel § 32 (9) anwendet:
\[ 2\eta(\omega)^4 = \eta\left(\frac{\omega}{2}\right)^2\vartheta_{00}\vartheta_{10}. \]

Multipliziert man beiderseits mit $\vartheta_{01}$, so folgt
nach (2):
\weq{7}{\eta\left(\frac{\omega}{2}\right)^2 = \vartheta_{01}\eta(\omega)}
und durch Verwandlung von $\omega$ in $\omega+1$ [(4) und § 31,
(6)]:
\weq{8}{e^{-\frac{\pi i}{12}}\eta\left(\frac{\omega+1}{2}\right)^2 = \vartheta_{00}\eta(\omega).}
